% Lecture 11: gRPC, Protocol Buffers, HTTP/2 and HTTP/3. Compile: latexmk -xelatex lecture11.tex
\documentclass[10pt,aspectratio=169,t]{beamer}
\usetheme{upbminimal}

\course{Mobile and Embedded Computing}
\decklabel{Lecture 11}
\title{gRPC, Protocol Buffers, HTTP/2 and HTTP/3}
\subtitle{Contract-first calls, streaming, and the transport underneath them}
\author{Dinu-Ștefan Rusu}
\institute{FILS, Universitatea Politehnica București}
\date{Spring 2027}

\begin{document}

\titleframe

\objectivesframe{
  \cell[\faServer]{Design the contract}{Write a .proto file and generate Go and Dart stubs with protoc, from one source of truth.}
  \cell[\faCode]{Call all four shapes}{Unary, server streaming, client streaming and bidirectional, and know which one a screen actually needs.}
  \cell[\faBolt]{Reason about the wire}{Why protobuf is smaller than JSON, how a schema evolves safely, and what HTTP/2 buys for free.}
  \cell[\faClock]{Survive the network}{Deadlines, cancellation and the radio tail: the habits that keep a mobile client polite.}
}

\begin{frame}[eyebrow=Overview]{What gRPC actually is}
  \vskip 3mm
  \begin{grid}[4]
    \cell[\faServer]{Contract first}{gRPC is a Remote Procedure Call (RPC) framework: one .proto file is the source of truth for every client and the server.}
    \cell[\faCode]{Code generation}{protoc turns messages and services into typed classes, in every language you target, from the same file.}
    \cell[\faBolt]{Binary and HTTP/2}{Compact binary frames, multiplexed on one warm connection instead of text on many short-lived ones.}
    \cell[\faIcon{sync-alt}]{Streaming built in}{Four call shapes are part of the contract itself, not a library bolted on afterward.}
  \end{grid}
  \note{Set the frame for the whole lecture: gRPC bundles a contract, a codec and a transport. Ask which of these four ideas the room has already met without the name attached to it.}
\end{frame}

\divider{Part 1 · Contracts}{Designing the call before writing code}[RPC versus REST, and the .proto contract]

\begin{frame}[fragile,eyebrow=Contracts]{RPC models verbs, REST models nouns}
  \begin{columns}[T]
    \begin{column}{0.53\textwidth}
      \begin{itemize}
        \item \textbf{REST (Representational State Transfer) models nouns.} \muted{You name resources and use a few fixed verbs on them, a good fit for a public API that strangers discover from a URL.}
        \item \textbf{RPC models verbs.} \muted{You name the operation your app needs, such as GetDashboard, and call it like a local function. Client and server agree on a list of procedures, not a URL space.}
        \item \textbf{On a phone the difference is round trips.} \muted{A resource-shaped screen often needs three or four sequential requests, and each one costs real time on cellular before any data moves.}
      \end{itemize}
    \end{column}
    \begin{column}{0.44\textwidth}
\begin{shell}
# REST: four requests
GET /users/42
GET /users/42/devices
GET /devices/9/readings
GET /devices/9/alerts

# RPC: one call
GetDashboard(user_id: 42)
  -> Dashboard
\end{shell}
      \begin{takeaway}{Design the call around the screen.}
        Not around your database tables. The round trip is the expensive part.
      \end{takeaway}
    \end{column}
  \end{columns}
  \note{Ask the room how many endpoints their ADMD project needed just to render one screen. That number is the argument for RPC before protobuf is even mentioned.}
\end{frame}

\begin{frame}[eyebrow=Contracts]{Choosing between REST and gRPC}
  \vskip 1mm
  \begin{center}
  {\usebeamerfont{small}
  \begin{tabular}{@{}p{32mm}p{40mm}p{40mm}@{}}
    \toprule
    \thead{Question} & \thead{Favor REST} & \thead{Favor gRPC} \\
    \midrule
    Who calls it? & Strangers, from a URL & You control both ends \\
    Needs streaming? & Rarely & Built into the contract \\
    Browser calls it? & Yes, plain HTTP & Only via gRPC-Web \\
    Payload format? & JSON, readable & Binary, needs tooling \\
    Latency budget? & One request is fine & Every round trip counts \\
    \bottomrule
  \end{tabular}}
  \end{center}
  \vskip 2mm
  \begin{takeaway}{Neither one is modern.}
    gRPC wins where you own both ends of the call. REST wins where you do not.
  \end{takeaway}
  \note{This table is the decision your team makes for project requirement four. Point out that a public, discoverable API is the one case where REST is still the right default.}
\end{frame}

\begin{frame}[fragile,eyebrow=Contracts]{Contract first: the .proto file}
  \begin{columns}[T]
    \begin{column}{0.62\textwidth}
\begin{codeblock}
syntax = "proto3";
package telemetry.v1;
message Reading {
  string device_id        = 1;
  double temp_c           = 2;
  int64  taken_at_unix_ms = 3;
}
service Telemetry {
  rpc SubmitReading(Reading) returns (SubmitReadingResponse);
  rpc WatchReadings(WatchReadingsRequest) returns (stream Reading);
}
\end{codeblock}
    \end{column}
    \begin{column}{0.35\textwidth}
      \lead{One file, three teams.}{Backend, Android and iOS build from it in parallel. Field numbers are the wire format: names never travel.}
      \muted{Enums need a zero: an absent enum decodes as 0, so 0 must mean UNSPECIFIED.}
    \end{column}
  \end{columns}
  \note{Write this live if there is time. The thing students miss: proto3 enums must have a zero value meaning unspecified, because an absent enum always decodes as 0. Versioning the package, not the file, is worth saying out loud too.}
\end{frame}

\begin{frame}[fragile,eyebrow=Contracts]{Generating stubs with protoc}
  \vskip 1mm
  Run protoc once per language, pointed at the same .proto file. Commit the generated code and regenerate on change, or run protoc in CI. Pick one and write it down.
  \vskip 3mm
\begin{shell}
# Go: two plugins, messages and service
protoc -Iproto \
  --go_out=gen      --go_opt=paths=source_relative \
  --go-grpc_out=gen --go-grpc_opt=paths=source_relative \
  proto/telemetry/v1/telemetry.proto

# Dart: the plugin is a pub global executable
dart pub global activate protoc_plugin
protoc -Iproto --dart_out=grpc:lib/gen \
  proto/telemetry/v1/telemetry.proto
\end{shell}
  \note{Live demo works well here: run these two commands and show the generated files appearing under gen/ and lib/gen. Both plugins read the exact same .proto file.}
\end{frame}

\begin{frame}[eyebrow=Contracts]{What the generator produces}
  \vskip 2mm
  \begin{itemize}
    \item \textbf{Typed message classes.} \muted{Reading, SubmitReadingResponse and their siblings, as real classes with typed getters, equality and a binary codec already written.}
    \item \textbf{A client stub.} \muted{TelemetryClient in Dart: the object your app calls as if it were a local class.}
    \item \textbf{A server skeleton.} \muted{UnimplementedTelemetryServer in Go: embed it, then implement one method at a time.}
    \item \textbf{Generated code is build output.} \muted{Never edit it by hand. The pubspec needs the protobuf and grpc packages to compile the result.}
  \end{itemize}
  \begin{takeaway}{The compiler writes the parsing code,}
    which is the code you would otherwise get subtly wrong.
  \end{takeaway}
  \note{Ask what happens if two developers hand-edit the same generated file differently. The answer is why generated code stays out of manual review.}
\end{frame}

\divider{Part 2 · Implementing}{A Go server and a generated Dart client}[Unary calls, then a stream of sensor readings]

\begin{frame}[fragile,eyebrow=Implementing]{A Go server: the unary call}
  \begin{columns}[T]
    \begin{column}{0.62\textwidth}
\begin{golang}
func (s *server) SubmitReading(ctx context.Context,
    req *pb.Reading) (*pb.SubmitReadingResponse, error) {
    id, err := s.store.Save(ctx, req)
    if err != nil {
        return nil, status.Error(codes.Internal, "save failed")
    }
    return &pb.SubmitReadingResponse{StoredId: id}, nil
}
\end{golang}
    \end{column}
    \begin{column}{0.35\textwidth}
      \lead{ctx is the deadline.}{Canceled when the phone gives up; the query stops here too.}
      \muted{Errors are codes, not strings: codes.InvalidArgument, Internal, and friends.}
    \end{column}
  \end{columns}
  \note{The server struct embeds pb.UnimplementedTelemetryServer so new RPCs keep it compiling. This is the same Go from Lecture 10, now speaking a generated interface. main() wires it up: net.Listen, grpc.NewServer, RegisterTelemetryServer, Serve.}
\end{frame}

\begin{frame}[eyebrow=Implementing]{Errors are codes, not strings}
  \vskip 1mm
  \begin{center}
  {\usebeamerfont{small}
  \begin{tabular}{@{}lp{85mm}@{}}
    \toprule
    \thead{Code} & \thead{Meaning} \\
    \midrule
    InvalidArgument   & The client sent something the server rejects outright \\
    NotFound          & No resource matches the request \\
    Unauthenticated   & No valid credentials were presented \\
    Unavailable       & The server is temporarily down, safe to retry \\
    DeadlineExceeded  & The call's deadline passed before a response arrived \\
    \bottomrule
  \end{tabular}}
  \end{center}
  \vskip 2mm
  \begin{takeaway}{The client branches on \code{e.code},}
    not on a parsed string. That is the whole point of a status code.
  \end{takeaway}
  \note{These codes are the same list on every language's generated client. A Dart catch block and a Go status.Error use the same vocabulary. PermissionDenied is worth a mention too: valid credentials, disallowed call.}
\end{frame}

\begin{frame}[fragile,eyebrow=Implementing]{A Dart client: the unary call}
  \begin{columns}[T]
    \begin{column}{0.6\textwidth}
\begin{dart}
final channel = ClientChannel('api.example.com', port: 443,
    options: const ChannelOptions(
        credentials: ChannelCredentials.secure()));
final client = TelemetryClient(channel);
final res = await client.submitReading(
  Reading()..deviceId = id..tempC = 21.5,
  options: CallOptions(timeout: const Duration(seconds: 2)),
);
debugPrint(res.storedId);
\end{dart}
    \end{column}
    \begin{column}{0.37\textwidth}
      {\usebeamerfont{small}
      \begin{itemize}
        \item \textbf{It looks like a local call.} \muted{That is the point of RPC, and the danger: it is still the network.}
        \item \textbf{Always set a deadline.} \muted{Without one, a stalled call holds a UI state forever.}
      \end{itemize}}
    \end{column}
  \end{columns}
  \note{Ask what happens to the UI if this call never sets a timeout and the server never responds. Someone will say it hangs forever, which is exactly right. A GrpcError catch block branches on e.code the same way the Go server used status codes.}
\end{frame}

\begin{frame}[fragile,eyebrow=Implementing]{Server streaming: the Go side}
  \begin{columns}[T]
    \begin{column}{0.58\textwidth}
\begin{golang}
func (s *server) WatchReadings(
    req *pb.WatchReadingsRequest,
    stream pb.Telemetry_WatchReadingsServer,
) error {
    ctx := stream.Context()
    for r := range s.store.Sub(ctx, req.DeviceId) {
        if err := stream.Send(r); err != nil {
            return err
        }
    }
    return nil
}
\end{golang}
    \end{column}
    \begin{column}{0.39\textwidth}
      \lead{One method, many responses.}{WatchReadings is declared once, as returns (stream Reading). The handler loops and sends until the client goes away.}
      \muted{ctx is still the deadline or the client's cancel. stream.Send fails once nobody is listening, and returning ends the stream cleanly.}
    \end{column}
  \end{columns}
  \note{This is the shape Lab 6 uses for live sensor readings, and the shape project requirement four needs if a team picks gRPC over FFI.}
\end{frame}

\begin{frame}[fragile,eyebrow=Implementing]{Server streaming: the Dart side}
  \begin{columns}[T]
    \begin{column}{0.55\textwidth}
\begin{dart}
final call = client.watchReadings(
  WatchReadingsRequest()..deviceId = id,
  options: CallOptions(timeout: const Duration(seconds: 5)));
final sub = call.listen(_onReading, onError: _onError);
@override
void dispose() {
  sub.cancel(); channel.shutdown();
  super.dispose();
}
\end{dart}
    \end{column}
    \begin{column}{0.42\textwidth}
      \lead{Cancel in dispose().}{An open stream keeps the radio and the server busy for a screen the user already left.}
      \muted{The deadline still applies: a stream silent that long fails too.}
    \end{column}
  \end{columns}
  \begin{takeaway}{Listen, cancel, and set a deadline:}
    the three habits Lab 6 grades directly.
  \end{takeaway}
  \note{Tie this straight to Lab 6's acceptance criteria: canceling stops the stream, and a short deadline against a server that sleeps three seconds fails visibly. call.listen returns a subscription like any other Dart Stream.}
\end{frame}

\begin{frame}[fragile,eyebrow=Implementing]{Client streaming: draining an offline queue}
  \begin{columns}[T]
    \begin{column}{0.47\textwidth}
      Many requests, one response: the shape of an outbox drain from Lecture 9. Queue rows offline, then stream them upstream for one acknowledgement.
    \end{column}
    \begin{column}{0.5\textwidth}
\begin{dart}
// service Telemetry adds:
// rpc UploadReadings(stream Reading) returns (Ack);
Stream<Reading> _drain() async* {
  for (final r in await outbox.rows()) {
    yield r.toReading();
  }
}
final sum = await client.uploadReadings(_drain());
debugPrint('stored: ${sum.accepted}');
\end{dart}
    \end{column}
  \end{columns}
  \note{Client streaming is the shape offline sync from Lecture 9 actually needs: drain the outbox in order and get one acknowledgement, not one round trip per row.}
\end{frame}

\divider{Part 3 · The wire}{What actually travels, and how it changes}[Protobuf's size on the wire, and evolving a schema safely]

\begin{frame}[fragile,eyebrow=Wire format]{Why protobuf is smaller, and by how much}
  \begin{columns}[T]
    \begin{column}{0.55\textwidth}
      {\usebeamerfont{small}
      \begin{itemize}
        \item \textbf{It is binary.} \muted{No quotes, commas or braces, and numbers are not rendered as text.}
        \item \textbf{Fields are numbered, not named.} \muted{A tag byte packs the field number and its wire type, so a field name never travels.}
        \item \textbf{Integers are varints.} \muted{Seven bits of value per byte, so small numbers cost one byte instead of four.}
        \item \textbf{Defaults are free.} \muted{Nothing is sent for a field left at its default value.}
      \end{itemize}}
    \end{column}
    \begin{column}{0.47\textwidth}
\begin{codeblock}[basicstyle=\ttfamily\fontsize{7}{9}\selectfont\color{upbInk}]
# JSON: 38 bytes
{"device_id":"esp32-a1","temp_c":21.5}
# protobuf: 19 bytes
0A 08 65 73 70 33 32 2D 61 31
11 00 00 00 00 00 80 35 40
# varint: 300 -> AC 02 (2 bytes)
\end{codeblock}
    \end{column}
  \end{columns}
  \begin{takeaway}{Quote the mechanism, not a multiplier.}
    Protobuf is smaller, often by a large factor, but the ratio depends on your schema and data. Measure yours.
  \end{takeaway}
  \note{This slide replaces the temptation to quote a multiplier. Have students serialize one real message both ways and print the byte counts themselves.}
\end{frame}

\begin{frame}[fragile,eyebrow=Wire format]{Schema evolution without breaking old apps}
  \begin{columns}[T]
    \begin{column}{0.55\textwidth}
\begin{codeblock}
message Reading {
  string device_id        = 1;
  double temp_c           = 2;
  int64  taken_at_unix_ms = 3;
  double humidity         = 5;  // new
  reserved 4;                   // retired
}
\end{codeblock}
    \end{column}
    \begin{column}{0.42\textwidth}
      {\usebeamerfont{small}
      \lead{Forward: old app, new server.}{The parser skips unknown tag 5 by its wire type. It does not crash.}
      \muted{Backward: a missing field decodes to its default. Never renumber, reuse, or retype a field.}}
    \end{column}
  \end{columns}
  \begin{takeaway}{Users do not update.}
    A year-old build is still on the network, and the wire format is the compatibility contract.
  \end{takeaway}
  \note{Ask what happens when an old build, still on someone's phone from last year, receives a message with a field it has never heard of. Nothing happens: it skips the bytes.}
\end{frame}

\begin{frame}[eyebrow=Call shapes]{Four call shapes, not one}
  {\usebeamerfont{small}\renewcommand\arraystretch{1.0}
  \begin{tabular}{@{}lp{42mm}p{50mm}@{}}
    \toprule
    \thead{Shape} & \thead{Signature} & \thead{When} \\
    \midrule
    Unary & \ttfamily rpc Get(Req) returns (Res) & One request, one response: most screens \\
    Server streaming & \ttfamily returns (stream Reading) & One request, many responses: live feeds \\
    Client streaming & \ttfamily (stream Reading) returns (Ack) & Many requests, one response: draining a queue \\
    Bidirectional & \ttfamily (stream In) returns (stream Out) & Both directions at once: chat, presence \\
    \bottomrule
  \end{tabular}}
  \begin{takeaway}{The shape is part of the contract.}
    Changing a method from unary to streaming changes the generated signature on every client.
  \end{takeaway}
  \note{Client streaming is the one students under-use: draining an offline queue is exactly this shape, one call instead of one round trip per row.}
\end{frame}

\divider{Part 4 · Transport}{HTTP/2, HTTP/3, and what a phone pays for}[One warm connection, and what happens when the network changes]

\begin{frame}[eyebrow=Transport]{HTTP/2: one connection, many streams}
  {\usebeamerfont{small}
  \begin{columns}[T]
    \begin{column}{0.55\textwidth}
      \begin{itemize}
        \item \textbf{HTTP/2 multiplexes calls over one connection.} \muted{HTTP/1.1 allowed one request at a time; browsers hid that with six connections. gRPC rides on this: one call per stream.}
        \item \textbf{What it does not fix: TCP delivers in order.} \muted{One lost packet stalls every stream on that connection until it is retransmitted.}
      \end{itemize}
    \end{column}
    \begin{column}{0.42\textwidth}
      \begin{itemize}
        \item \textbf{One TCP + TLS connection.} \muted{Handshake once, reuse it for every call.}
        \item \textbf{HPACK, HTTP/2's header compression.} \muted{A repeated header costs a table index, not the bytes.}
        \item \textbf{Streams are independent.} \muted{A slow call does not block a fast one.}
      \end{itemize}
    \end{column}
  \end{columns}}
  \begin{takeaway}{One connection, kept warm.}
    That is most of gRPC's mobile advantage, before a single byte of protobuf is counted.
  \end{takeaway}
  \note{The head-of-line blocking TCP still causes is the setup for the next slide: QUIC is the fix HTTP/2 could not make on its own. HPACK is HTTP/2's header compression.}
\end{frame}

\begin{frame}[eyebrow=Transport]{HTTP/3 and QUIC: the mobile case}
  {\usebeamerfont{small}
  \begin{columns}[T]
    \begin{column}{0.55\textwidth}
      \begin{itemize}
        \item \textbf{QUIC moves onto UDP and implements streams itself.} \muted{A lost packet now stalls only that stream, so HTTP/2's head-of-line blocking goes away.}
        \item \textbf{Transport Layer Security (TLS) 1.3 is folded into the handshake.} \muted{A new connection is one round trip; a resumed one can be zero.}
        \item \textbf{Connection migration matters here.} \muted{A QUIC connection has an id, not an IP and port, so it survives Wi-Fi to 5G.}
      \end{itemize}
    \end{column}
    \begin{column}{0.42\textwidth}
      \begin{itemize}
        \item \textbf{Where gRPC stands.} \muted{The wire protocol is specified over HTTP/2.}
        \item \textbf{In practice.} \muted{HTTP/3 usually reaches you at the edge: a load balancer speaks it to the device, HTTP/2 to your service.}
        \item \textbf{Treat it as deployment,} \muted{not a code change.}
      \end{itemize}
    \end{column}
  \end{columns}}
  \begin{takeaway}{Handshakes and reconnects are the mobile cost.}
    QUIC addresses both, which is why it belongs in the same argument as the radio tail.
  \end{takeaway}
  \note{Emphasize connection migration: it is the one QUIC feature with no equivalent fix in HTTP/2, because TCP's identity is tied to the IP and port.}
\end{frame}

\begin{frame}[eyebrow=Numbers]{One connection, four shapes, zero round trips}
  \vskip 4mm
  \begin{grid}[3]
    \bignum{1}{TCP connection, reused for every call}
    \bignum{4}{call shapes defined by the .proto, not by you}
    \bignum{0-RTT}{round-trip time on a resumed QUIC connection}
  \end{grid}
  \note{These numbers are structural, not benchmarks: they hold regardless of network or hardware, which is why they are safe to state without a citation.}
\end{frame}

\begin{frame}[eyebrow=On a phone]{Deadlines, cancellation, and the radio tail}
  {\usebeamerfont{small}
  \begin{columns}[T]
    \begin{column}{0.55\textwidth}
      \begin{itemize}
        \item \textbf{A deadline is not a client-side timeout.} \muted{It travels with the call and propagates through every service behind it. When it expires, work stops everywhere.}
        \item \textbf{The radio has a tail.} \muted{It stays high-power for seconds after a transmission. One burst allows it to sleep; several chatty requests keep it awake.}
      \end{itemize}
    \end{column}
    \begin{column}{0.42\textwidth}
      \begin{hint}[iOS caveat]
        A suspended app gets no CPU and its sockets close, so a long-lived stream does not survive backgrounding. Reconnect on resume; use a silent push or a background task to be woken.
      \end{hint}
    \end{column}
  \end{columns}}
  \begin{takeaway}{Batch, set a deadline, cancel on dispose.}
    Three habits that show up as battery life in a review.
  \end{takeaway}
  \note{This is where Lecture 7's background execution material and this lecture meet. A team whose project needs a live gRPC stream while backgrounded will find it does not survive on iOS. Keepalive pings are not free either: they hold the connection open and wake the radio, so tune the interval.}
\end{frame}

\divider{Part 5 · Reaching everywhere}{The browser, and the project}[gRPC-Web, Connect, and where Lab 6 picks this up]

\begin{frame}[fragile,eyebrow=gRPC]{gRPC-Web and Connect: the browser problem}
  \begin{columns}[T]
    \begin{column}{0.55\textwidth}
      {\usebeamerfont{small}
      \begin{itemize}
        \item \textbf{A browser cannot speak gRPC.} \muted{fetch and XHR give no control over HTTP/2 frames and no way to read trailers, where gRPC puts the call's status.}
        \item \textbf{gRPC-Web} \muted{is a different wire encoding, over HTTP/1.1, for unary and server streaming only. It needs a translating proxy.}
        \item \textbf{Connect} \muted{is the pragmatic successor: one Go handler serves all three protocols on the same route.}
      \end{itemize}}
    \end{column}
    \begin{column}{0.42\textwidth}
\begin{golang}[basicstyle=\ttfamily\fontsize{7}{9}\selectfont\color{upbInk}]
mux := http.NewServeMux()
path, h := connect.NewHandler(
    &server{})
mux.Handle(path, h)
http.ListenAndServe(":8080", mux)
\end{golang}
      \muted{A Flutter web build cannot open a raw gRPC connection.}
    \end{column}
  \end{columns}
  \begin{takeaway}{The boundary has a shape here too:}
    what the browser sandbox will let you send is part of your architecture.
  \end{takeaway}
  \note{If a team's project ships a Flutter web build, this is the slide that tells them why a raw gRPC channel will not compile for that target.}
\end{frame}

\begin{frame}[eyebrow=Project]{Where this lands in Lab 6 and the project}
  \vskip 3mm
  \begin{grid}[2]
    \cell[\faServer]{Lab 6}{One .proto with a unary call and a server-streaming call, a Go server, and a generated Dart client. The stream updates live, canceling stops it, and a two-second deadline fails a call against a server that sleeps three seconds.}
    \cell[\faIcon{tasks}]{The project}{Requirement four accepts gRPC as one boundary choice: a Go server with a generated Dart client, alongside FFI or a platform channel from Lecture 12.}
  \end{grid}
  \begin{takeaway}{Everything on this slide, you now have the vocabulary for.}
    The exact commands and acceptance criteria are in the Lab 6 handout.
  \end{takeaway}
  \note{Preview Lab 6 concretely: a .proto with one unary and one server-streaming call, generated stubs, a running server, a Flutter client, a deadline and a cancel.}
\end{frame}

\begin{frame}[eyebrow=Recap]{Three things to remember}
  \vskip 2mm
  \begin{enumerate}
    \item A .proto is a contract, not a suggestion. \muted{Field numbers are the wire format; never renumber or reuse one.}
    \item The shape of a call is part of the contract. \muted{Unary, server streaming, client streaming and bidirectional are four different generated signatures.}
    \item Deadlines and cancellation make gRPC polite on a phone. \muted{Set one, cancel in dispose, and mind the radio tail.}
  \end{enumerate}
  \vskip 5mm
  \kv{Further reading}{\link{https://grpc.io/docs/}{grpc.io/docs} \textperiodcentered\ \link{https://protobuf.dev/programming-guides/encoding/}{protobuf.dev/encoding} \textperiodcentered\ \link{https://www.rfc-editor.org/rfc/rfc9114}{RFC 9114 (HTTP/3)}}
  \note{Close by connecting back to the objectives slide: contract, wire, transport, and the mobile habits are the same four ideas from the start of the lecture.}
\end{frame}

\closingframe{Questions?}{dinu\_stefan.rusu@upb.ro · Microsoft Teams: course channel}

\end{document}
